\documentclass{article}
\usepackage{fuzz}
\makeatletter
\newenvironment{theorem}[1]{\medskip\noindent\textbf{Theorem (#1).}\@nameuse{zed}}{\@nameuse{endzed}}
\makeatother

\begin{document}

\title{A Formal Specification of \texttt{model-provider} (current state: faithful to \texttt{model-provider.sh} as built)}
\author{Punt Labs}
\maketitle

\section{Purpose and scope}

This document models \texttt{model-provider.sh} \emph{as it is implemented
today}. It is the faithful companion to \texttt{model-provider.target.spec.tex},
which models the more robust design we are building toward. Where the two
differ, this one describes the code; the target one describes the intent. Each
schema here was checked against the script, and every divergence a peer review
found between an earlier, idealised draft and the real code is resolved \emph{in
favour of the code} below.

The current tool takes one operator selection --- a model served by the
\textbf{Bifrost} gateway (the only gateway it implements) --- and configures
four AI coding harnesses (Claude Code, codex, pi, opencode). The guarantee it
actually delivers is weaker and more specific than ``every request serves the
selection''. First, the \emph{main} request of every harness serves the
selection. Second, Claude Code's \emph{subagent} (tier-alias) and \emph{small}
requests are routed to a \emph{different, same-provider-family} id resolved from
the live catalog --- the highest-version model of that tier in the selection's
provider family --- not to the selection itself; a haiku subagent is meant to
run haiku, not to be force-upgraded onto the selection. Third, \texttt{verify}
probes \emph{only} Claude Code's Anthropic path against Bifrost; on the direct
route it is a successful no-op, not a failure.

Preconditions are written explicitly so TTF partition analysis can derive test
cases from each operation's input domain.

\section{Given sets and enumerations}

A model id is an opaque catalog token (e.g.
\texttt{openrouter/anthropic/claude-opus-4.8}). A provider is the upstream
Bifrost routes a model to (e.g. \texttt{openrouter}, \texttt{anthropic},
\texttt{bedrock}) --- in a model id it is the segment before the first slash. A
canonical id is the base model name Claude Code derives from a model id (e.g.
\texttt{claude-opus-4-8}).

\begin{zed}
  [ModelId, Provider, CanonId]
\end{zed}

The four harnesses.

\begin{zed}
  Harness ::= claudecode | codex | pi | opencode
\end{zed}

Each harness issues requests of distinct kinds. \emph{main} is the primary
loop; \emph{subagent} is a Claude Code Task-tool spawn (resolved through tier
alias configuration); \emph{small} is a background / side model (Claude Code's
small-fast model, opencode's title model). Side-requests are first-class: they
are exactly the paths that fail silently when only the main model is set.

\begin{zed}
  ReqKind ::= main | subagent | small
\end{zed}

The route the tool can be in. The current tool implements exactly two: routed
through Bifrost, or ``direct'' (vendor defaults, no gateway). There is no
gateway selection --- Bifrost is hardcoded --- so unlike the target spec there
is no \emph{set} of gateways here, only this route flag.

\begin{zed}
  Route ::= bifrost | direct
\end{zed}

Authentication modes a Bifrost route may use: a per-host virtual key
($vkauth$), or the operator's own subscription OAuth passed through Bifrost
($subauth$).

\begin{zed}
  AuthMode ::= vkauth | subauth
\end{zed}

Booleans. fuzz's mathematical toolkit has no built-in boolean type, so the
specification defines one; $yes$/$no$ are its values.

\begin{zed}
  Bool ::= yes | no
\end{zed}

\section{Harness properties}

Which harnesses \emph{canonicalize} the configured id before sending it (strip
the provider prefix, rewrite the version), rather than routing it verbatim.
Only Claude Code does; this is why it needs an override map to send a
dotted/prefixed id unaltered on the wire.

\begin{axdef}
  canonicalizes : \power Harness
\where
  canonicalizes = \{ claudecode \}
\end{axdef}

The request kinds each harness actually issues.

\begin{axdef}
  issues : Harness \fun \power ReqKind
\where
  issues = \{ claudecode \mapsto \{ main, subagent, small \}, \\
  \t1         codex \mapsto \{ main \}, \\
  \t1         pi \mapsto \{ main \}, \\
  \t1         opencode \mapsto \{ main, small \} \}
\end{axdef}

\section{The Bifrost catalog}

The tool reads one live catalog from Bifrost: per model id, the provider it
routes to. $catalog$ is the set of ids Bifrost knows; the operator may select
any of them (validation against this set is exactly what \texttt{use} does
without \texttt{--force}). $provOf$ gives the upstream provider of a model id.

\begin{axdef}
  catalog : \power ModelId \\
  provOf : ModelId \pfun Provider
\where
  \dom provOf = catalog \\
  catalog \neq \emptyset
\end{axdef}

\section{Tier resolution --- the crux of the current design}

For Claude Code, the subagent and small request kinds are \emph{not} served by
the selection. \texttt{claude\_tier} resolves them, per tier, to the
highest-version catalog id of that tier \emph{in the selection's provider
family}. $tierId$ models this: given the selection, it yields the id a tier
slot is set to. It is a same-family id --- its provider equals the selection's
provider --- but it is generally a \emph{different model}. On a non-Claude
selection (no Claude tier exists in that family) the slot is omitted; we model
that with the selection's own id standing in only to keep the slot total, but
the fidelity property we assert is provider-equality, not id-equality.

\begin{axdef}
  tierId : ModelId \fun ModelId
\where
  \forall m : ModelId @ \\
  \t1 m \in catalog \implies \\
  \t2 ( tierId~m \in catalog \land provOf~(tierId~m) = provOf~m )
\end{axdef}

The id that reaches the wire for harness $h$, request kind $k$, under selection
$m$, is written inline where it is used, as the conditional expression
\[ \IF (h = claudecode \land k \neq main) \THEN tierId~m \ELSE m. \]
The main kind of every harness, and every kind of a non-canonicalizing harness
that has only side-slots pinned to the selection, carries $m$ verbatim; Claude
Code's subagent and small kinds carry $tierId~m$ --- a same-family,
generally-different id. This is the honest current-state wire target: it does
\emph{not} claim the selection reaches every request. It is inlined rather than
packaged as a standalone function because a total function over the
$Harness \cross ReqKind \cross ModelId$ product is not finitely enumerable by
the model checker, whereas the inline form ranges only over each harness's
finite $issues~h$.

\section{Configuration state}

$sel$ is the selection; $route$ records bifrost (routed) or direct; $auth$ the
auth mode. $wire$ gives, per harness, per issued request kind, the id that will
actually travel. $overrides$ records, per harness, whether a canonicalization
override has been written. The selection must be a catalog id \emph{unless} it
was forced; $forced$ records that, so a \texttt{use --force} of an off-catalog
id is representable (the real tool writes it to the block).

\begin{schema}{Config}
  sel : ModelId \\
  route : Route \\
  auth : AuthMode \\
  forced : Bool \\
  wire : Harness \pfun (ReqKind \pfun ModelId) \\
  overrides : \power Harness
\where
  forced = no \implies sel \in catalog
\end{schema}

\section{Initial state}

A fresh repo reads \texttt{BIFROST\_ROUTE} with a default of \texttt{bifrost},
but with no selection and no managed wire --- \texttt{status} reports
\texttt{route: bifrost} with \texttt{model: <unset>}. We model the never-used
state as direct with an empty wire, matching what the managed block actually
contains (no \texttt{ANTHROPIC\_MODEL} line) before the first \texttt{use};
the route-default read-back is a status concern, not a wire state.

\begin{schema}{Init}
  Config'
\where
  route' = direct \\
  auth' = vkauth \\
  forced' = no \\
  wire' = \{ h : \{ claudecode, codex, pi, opencode \} @ h \mapsto \emptyset \} \\
  overrides' = \emptyset
\end{schema}

\section{Well-formedness}

$wire$ is defined for exactly the four harnesses and, for each, exactly the
request kinds it issues; overrides are only ever recorded for canonicalizing
harnesses.

\begin{schema}{WellFormed}
  Config
\where
  \dom wire = \{ claudecode, codex, pi, opencode \} \\
  ( \forall h : \dom wire @ \dom (wire~h) = issues~h ) \\
  overrides \subseteq canonicalizes
\end{schema}

\section{The core guarantee (as built)}

\textbf{Routing correctness, current.} When routed through Bifrost, every
request every harness issues carries the inline wire target --- the
selection for main kinds, the same-family tier id for Claude Code's
side-requests --- and Claude Code has the override that stops its main id being
re-canonicalized.

\begin{schema}{RoutingCorrectCur}
  WellFormed
\where
  route = bifrost \implies \\
  \t1 ( ( \forall h : \dom wire @ \forall k : \dom (wire~h) @ \\
  \t2     (wire~h)~k = ( \IF (h = claudecode \land k \neq main) \THEN tierId~sel \ELSE sel ) ) \land \\
  \t1   canonicalizes \subseteq overrides )
\end{schema}

\textbf{Provider fidelity --- the real invariant.} The honest guarantee the
current code makes: every request, main or side, reaches the \emph{selection's
provider}, even when a side-request runs a different model in that family.
This is weaker than id-fidelity and is exactly what \texttt{claude\_tier}
guarantees by construction (it only ever picks within the selection's family).

\begin{schema}{ProviderFidelityCur}
  RoutingCorrectCur
\where
  route = bifrost \implies \\
  \t1 ( \forall h : \dom wire @ \forall k : \dom (wire~h) @ \\
  \t2   provOf~((wire~h)~k) = provOf~sel )
\end{schema}

\textbf{The forbidden failure.} The historical bug: a canonicalizing harness
sent the base id, or a side-request fell back to an unroutable built-in whose
provider is not the selection's. Named as the predicate the configuration makes
unreachable while routed --- the negation of the two guarantees above.

\begin{schema}{SilentMisrouteCur}
  WellFormed
\where
  route = bifrost \\
  ( ( \exists h : \dom wire @ \exists k : \dom (wire~h) @ \\
  \t1   provOf~((wire~h)~k) \neq provOf~sel ) \lor \\
  \t1 ( \exists h : canonicalizes @ h \notin overrides ) )
\end{schema}

\section{Shared effect: applying the Bifrost configuration}

The effect every selecting / re-applying operation shares: route through
Bifrost with selection $m$, rebuild $wire$ so each issued request kind carries
the inline wire target (the selection, or its same-family tier id for Claude
Code's side-requests), and record overrides for exactly the canonicalizing
harnesses. $fwd?$ carries the forced flag through to state.

\begin{schema}{ApplyBifrost}
  \Delta Config \\
  m? : ModelId \\
  fwd? : Bool
\where
  sel' = m? \\
  route' = bifrost \\
  auth' = auth \\
  forced' = fwd? \\
  \dom wire' = \{ claudecode, codex, pi, opencode \} \\
  ( \forall h : \dom wire' @ \\
  \t1 wire'~h = \{ k : issues~h @ \\
  \t2   k \mapsto ( \IF (h = claudecode \land k \neq main) \THEN tierId~m? \ELSE m? ) \} ) \\
  overrides' = canonicalizes
\end{schema}

\section{Operation: \texttt{use}}

Selecting a model. Input is a model id and a $force$ flag (no gateway input ---
Bifrost is hardcoded). The precondition domain partitions into: (a) valid --- id
in the live catalog; (b) forced --- $force? = yes$ bypasses the catalog check;
(c) refused --- id not in catalog and not forced. An id that strips to empty is
refused even under force, but that is a lexical guard on the raw argument, below
the model layer, so it is noted here in prose rather than as a state schema.

\begin{schema}{UseOk}
  ApplyBifrost \\
  force? : Bool
\where
  m? \in catalog \\
  fwd? = force? \\
  force? = no
\end{schema}

A forced selection: the id need not be in the catalog, and the forced flag is
recorded so $Config$ admits the off-catalog $sel$. This partition is
\emph{reachable} here (unlike the target spec's idealised version) precisely
because $Config$ carries $forced$.

\begin{schema}{UseForced}
  ApplyBifrost \\
  force? : Bool
\where
  force? = yes \\
  fwd? = yes
\end{schema}

The refusal: id unknown to the catalog and not forced. No state change.

\begin{schema}{UseRefused}
  \Xi Config \\
  m? : ModelId \\
  force? : Bool \\
  err! : \num
\where
  m? \notin catalog \\
  force? = no \\
  err! = 1
\end{schema}

\begin{zed}
  Use \defs (UseOk \lor UseForced) \lor UseRefused
\end{zed}

\section{Operation: \texttt{auth}}

Set the auth mode and \emph{re-apply the current route unchanged}. The real
tool passes the current route through (\texttt{write\_block "\$route" \ldots}):
from bifrost it re-applies bifrost; from direct it stays direct. The
\texttt{require\_model\_for\_bifrost} guard refuses only when the current route
is bifrost \emph{and} no selection has ever been made --- so \texttt{direct}
then \texttt{auth vk} must succeed and stay direct. We model the selection-made
precondition as $sel \in catalog \lor forced = yes$ (a real selection exists).

\begin{schema}{AuthOkBifrost}
  \Delta Config \\
  newauth? : AuthMode
\where
  route = bifrost \\
  ( sel \in catalog \lor forced = yes ) \\
  sel' = sel \\
  route' = bifrost \\
  auth' = newauth? \\
  forced' = forced \\
  \dom wire' = \{ claudecode, codex, pi, opencode \} \\
  ( \forall h : \dom wire' @ \\
  \t1 wire'~h = \{ k : issues~h @ \\
  \t2   k \mapsto ( \IF (h = claudecode \land k \neq main) \THEN tierId~sel \ELSE sel ) \} ) \\
  overrides' = canonicalizes
\end{schema}

On the direct route, \texttt{auth} changes only the recorded auth mode: the
direct branch of \texttt{write\_block} ignores the model and writes no wire, so
no \texttt{ANTHROPIC\_MODEL} is touched and the guard does not fire.

\begin{schema}{AuthOkDirect}
  \Delta Config \\
  newauth? : AuthMode
\where
  route = direct \\
  sel' = sel \\
  route' = direct \\
  auth' = newauth? \\
  forced' = forced \\
  wire' = wire \\
  overrides' = overrides
\end{schema}

The refusal: current route is bifrost but no selection has ever been made.

\begin{schema}{AuthRefusedNoModel}
  \Xi Config \\
  newauth? : AuthMode \\
  err! : \num
\where
  route = bifrost \\
  sel \notin catalog \\
  forced = no \\
  err! = 1
\end{schema}

\begin{zed}
  Auth \defs (AuthOkBifrost \lor AuthOkDirect) \lor AuthRefusedNoModel
\end{zed}

\section{Operations: \texttt{bifrost}, \texttt{direct}, \texttt{toggle}}

$GatewayOn$ (the \texttt{bifrost} command) re-applies the current selection on
Bifrost; it is refused with no selection (the same
\texttt{require\_model\_for\_bifrost} guard). $Direct$ withdraws the managed
wire and overrides, retaining $sel$/$auth$ so a later $GatewayOn$ re-applies
them. \texttt{toggle} flips the route, applying the same no-selection guard
when it would turn Bifrost on.

\begin{schema}{GatewayOn}
  \Delta Config
\where
  ( sel \in catalog \lor forced = yes ) \\
  sel' = sel \\
  route' = bifrost \\
  auth' = auth \\
  forced' = forced \\
  \dom wire' = \{ claudecode, codex, pi, opencode \} \\
  ( \forall h : \dom wire' @ \\
  \t1 wire'~h = \{ k : issues~h @ \\
  \t2   k \mapsto ( \IF (h = claudecode \land k \neq main) \THEN tierId~sel \ELSE sel ) \} ) \\
  overrides' = canonicalizes
\end{schema}

\begin{schema}{Direct}
  \Delta Config
\where
  route' = direct \\
  sel' = sel \\
  auth' = auth \\
  forced' = forced \\
  wire' = \{ h : \{ claudecode, codex, pi, opencode \} @ h \mapsto \emptyset \} \\
  overrides' = \emptyset
\end{schema}

\begin{schema}{Toggle}
  \Delta Config
\where
  sel' = sel \\
  auth' = auth \\
  forced' = forced \\
  ( route = direct \implies \\
  \t1 ( route' = bifrost \land ( sel \in catalog \lor forced = yes ) \land \\
  \t2   \dom wire' = \{ claudecode, codex, pi, opencode \} \land \\
  \t2   ( \forall h : \dom wire' @ \\
  \t3     wire'~h = \{ k : issues~h @ \\
  \t4       k \mapsto ( \IF (h = claudecode \land k \neq main) \THEN tierId~sel \ELSE sel ) \} ) \land \\
  \t2   overrides' = canonicalizes ) ) \\
  ( route = bifrost \implies \\
  \t1 ( route' = direct \land \\
  \t2   wire' = \{ h : \{ claudecode, codex, pi, opencode \} @ h \mapsto \emptyset \} \land \\
  \t2   overrides' = \emptyset ) )
\end{schema}

\section{Read-only operations: \texttt{status}, \texttt{list}, \texttt{verify}}

$Status$ reports the selection, route and auth. $List$ reports the catalog.
$Verify$ is scoped to what the real \texttt{verify\_model} does: it probes
\emph{only Claude Code's} Anthropic path against Bifrost. On the direct route it
is a successful no-op (the tool prints a note and returns success) --- this is
$VerifyDirectOk$, not a failure.

\begin{schema}{Status}
  \Xi Config \\
  repSel! : ModelId \\
  repRoute! : Route \\
  repAuth! : AuthMode
\where
  repSel! = sel \\
  repRoute! = route \\
  repAuth! = auth
\end{schema}

\begin{schema}{List}
  \Xi Config \\
  cat! : \power ModelId
\where
  cat! = catalog
\end{schema}

\texttt{verify} succeeds when, routed through Bifrost, Claude Code's \emph{main}
request carries the selection and the override is present. It deliberately does
\emph{not} check codex/pi/opencode (those are separate manual checks in the
test plan) and does not assert the selection on Claude Code's side-requests
(they carry tier ids). $WellFormed$ is included so a missing Claude Code main
slot cannot vacuously pass.

\begin{schema}{VerifyPass}
  \Xi Config \\
  WellFormed \\
  ok! : \num
\where
  route = bifrost \\
  (wire~claudecode)~main = sel \\
  claudecode \in overrides \\
  ok! = 1
\end{schema}

On the direct route, \texttt{verify} is a successful no-op: nothing routes
through Bifrost, so there is nothing to prove.

\begin{schema}{VerifyDirectOk}
  \Xi Config \\
  ok! : \num
\where
  route = direct \\
  ok! = 1
\end{schema}

The genuine failure: routed, but Claude Code's main request does not carry the
selection, or the override is missing.

\begin{schema}{VerifyFail}
  \Xi Config \\
  ok! : \num
\where
  route = bifrost \\
  ( (wire~claudecode)~main \neq sel \lor claudecode \notin overrides ) \\
  ok! = 0
\end{schema}

\begin{zed}
  Verify \defs (VerifyPass \lor VerifyDirectOk) \lor VerifyFail
\end{zed}

\section{The system invariant}

Whenever routed through Bifrost, the configuration is well-formed, every
request carries the inline wire target, and every request reaches the
selection's provider. This is the precise ``works as built'' property. It holds
by construction of the operations (each rebuilds $wire$ via the inline target); the
theorem below makes that machine-checkable.

\begin{schema}{SystemOKCur}
  WellFormed
\where
  route = bifrost \implies \\
  \t1 ( ( \forall h : \dom wire @ \forall k : \dom (wire~h) @ \\
  \t2     (wire~h)~k = ( \IF (h = claudecode \land k \neq main) \THEN tierId~sel \ELSE sel ) ) \land \\
  \t1   ( \forall h : \dom wire @ \forall k : \dom (wire~h) @ \\
  \t2     provOf~((wire~h)~k) = provOf~sel ) \land \\
  \t1   canonicalizes \subseteq overrides )
\end{schema}

A machine-checkable obligation: every state reached by the shared apply effect
satisfies the system invariant. $ApplyBifrost$'s postcondition builds a
well-formed $wire$ from the inline target with overrides for the canonicalizing
harnesses, so the invariant follows. Stating it as a theorem lets
\texttt{probcli -cbc\_assertions} actually check it, rather than relying on
prose.

\begin{theorem}{ApplyGivesSystemOKCur}
  \forall Config'; m? : ModelId; fwd? : Bool @ \\
  \t1 ( \exists Config @ ApplyBifrost ) \implies \\
  \t2 ( \theta Config' \in \{ SystemOKCur @ \theta Config \} )
\end{theorem}

\section{Per-harness on-disk configuration}

The abstract $wire$ states the guarantee; this section grounds it in each
harness's actual files and model slots, including sub-agent / extension slots,
as confirmed against the on-disk configs. A model slot holds a concrete id, or
$inherit$ (unset --- the harness falls back to its parent / main model).

\begin{zed}
  Slot ::= inherit | fixed \ldata ModelId \rdata
\end{zed}

A slot \emph{resolves to the selection} when it is $inherit$ or $fixed~sel$.
For Claude Code's tier/small slots the current code does \emph{not} target the
selection, so a second predicate, $resolvesFamily$, captures the real
guarantee: the slot is $inherit$, or a $fixed$ id in the selection's provider
family.

\begin{axdef}
  resolves : ModelId \fun \power Slot
\where
  \forall m : ModelId @ \\
  \t1 resolves~m = \{ inherit \} \cup \{ s : Slot | s = fixed~m \}
\end{axdef}

\begin{axdef}
  resolvesFamily : ModelId \fun \power Slot
\where
  \forall m : ModelId @ \\
  \t1 resolvesFamily~m = \{ inherit \} \cup \\
  \t2 \{ s : Slot; i : ModelId | s = fixed~i \land i \in catalog \land \\
  \t3   provOf~i = provOf~m @ s \}
\end{axdef}

\subsection{Claude Code}

Config lives in the managed block in $DIR$\texttt{/.envrc.local} and
\texttt{.claude/settings.local.json} (\texttt{modelOverrides}). Slots: $ccMain$
(\texttt{ANTHROPIC\_MODEL}); the four tier-alias slots
$ccOpus$/$ccSonnet$/$ccHaiku$/$ccFable$
(\texttt{ANTHROPIC\_DEFAULT\_\{OPUS,SONNET,HAIKU,FABLE\}\_MODEL}); and $ccSmall$
(\texttt{ANTHROPIC\_SMALL\_FAST\_MODEL}). $ccOverride$ records the
\texttt{modelOverrides} map (canonical-key and raw-identity entries).

\begin{schema}{ClaudeCodeConfig}
  ccMain, ccOpus, ccSonnet, ccHaiku, ccFable, ccSmall : Slot \\
  ccOverride : Bool
\end{schema}

Claude Code is correct, as built, when the \emph{main} slot resolves to the
selection and the override is present, and each tier/small slot resolves
\emph{within the selection's family} (it is set to a same-family tier id, or
left to inherit). This is the faithful \texttt{ClaudeCodeOK}: it does not demand
the tier slots equal the selection, because the code does not set them that way.

\begin{schema}{ClaudeCodeOKCur}
  ClaudeCodeConfig \\
  sel : ModelId
\where
  ccMain \in resolves~sel \\
  ccOpus \in resolvesFamily~sel \land ccSonnet \in resolvesFamily~sel \\
  ccHaiku \in resolvesFamily~sel \land ccFable \in resolvesFamily~sel \\
  ccSmall \in resolvesFamily~sel \\
  ccOverride = yes
\end{schema}

\subsection{codex}

Config lives in \texttt{\~{}/.codex/config.toml}: $cxMain$ (\texttt{model}) and
the sub-agent slot $cxSubagent$ (\texttt{[agents]}
\texttt{agent\_default\_subagent\_model}). A spawned sub-agent that omits a
model inherits $cxMain$; model-provider leaves it unset, so $inherit$ is both
default and safe.

\begin{schema}{CodexConfig}
  cxMain : Slot \\
  cxSubagent : Slot
\end{schema}

\begin{schema}{CodexOKCur}
  CodexConfig \\
  sel : ModelId
\where
  cxMain \in resolves~sel \\
  cxSubagent \in resolves~sel
\end{schema}

\subsection{pi}

Config lives in \texttt{\~{}/.pi/agent/settings.json} ($piMain$,
\texttt{defaultModel}); sub-agents are Markdown files whose omitted frontmatter
\texttt{model:} inherits the session model. $piSubagent$ models that slot;
$inherit$ is the safe default.

\begin{schema}{PiConfig}
  piMain : Slot \\
  piSubagent : Slot
\end{schema}

\begin{schema}{PiOKCur}
  PiConfig \\
  sel : ModelId
\where
  piMain \in resolves~sel \\
  piSubagent \in resolves~sel
\end{schema}

\subsection{opencode}

Config lives in \texttt{\~{}/.config/opencode/opencode.json}: $ocMain$
(\texttt{model}), $ocSmall$ (\texttt{small\_model}), and per-agent slots
$ocAgent$ (\texttt{agent.<name>.model}). The current tool pins $ocSmall$ to the
selection (\texttt{small\_model = openai/<model>}, proven routable), so for
opencode even the small slot resolves to the selection --- unlike Claude Code.
An agent that omits \texttt{model} inherits the session model.

\begin{schema}{OpencodeConfig}
  ocMain : Slot \\
  ocSmall : Slot \\
  ocAgent : Slot
\end{schema}

\begin{schema}{OpencodeOKCur}
  OpencodeConfig \\
  sel : ModelId
\where
  ocMain \in resolves~sel \\
  ocSmall \in resolves~sel \\
  ocAgent \in resolves~sel
\end{schema}

\subsection{The composed fleet invariant}

When routed, Claude Code is correct to \emph{family} fidelity on its side-slots
and \emph{id} fidelity on main; codex, pi and opencode are correct to id
fidelity. This is the file-level, sub-agent-inclusive restatement of
$SystemOKCur$.

\begin{schema}{Fleet}
  sel : ModelId \\
  route : Route \\
  ClaudeCodeConfig \\
  CodexConfig \\
  PiConfig \\
  OpencodeConfig
\end{schema}

\begin{schema}{FleetOKCur}
  Fleet
\where
  route = bifrost \implies \\
  \t1 ( ClaudeCodeOKCur \land CodexOKCur \land PiOKCur \land OpencodeOKCur )
\end{schema}

\section{What this spec does not claim}

For honesty, the current-state limitations modelled or noted above are as
follows. Only Bifrost is implemented; there is no gateway selection.
Entitlement is not enforced: the catalog domain is the whole selectable set, so
a catalogued-but-unentitled id (the Bedrock Anthropic-models gap) can be
selected and routed --- the target spec closes this. \texttt{verify} proves
only Claude Code's main Anthropic path; the other harnesses and side-requests
are verified by the separate manual procedure in the test plan. Finally,
\texttt{direct} leaves some opencode provider/model residue in place (a
documented cleanup gap); $Direct$ here idealises that to an empty wire --- the
one place this spec is slightly kinder than the code, called out rather than
hidden.

\section{Two real-world constraints: compatibility and credential terms}

Two real-world facts bound correct operation; neither is enforced by the
current tool, so both are documented gaps.

\textbf{Compatibility (support, not permission).} Anthropic does \emph{not}
support routing Claude Code to non-Claude models. This is a support and
compatibility statement, not a terms violation: with a non-Anthropic selection
on the Claude Code path, tool use, extended thinking, and prompt caching are
unvalidated and degrade per model, so routing fidelity is necessary but not
sufficient. Separately, two Claude Code defaults --- MCP tool search and
fine-grained tool streaming --- silently turn off whenever Claude Code is routed
through the gateway (the \texttt{bifrost} route), regardless of model; both are
recoverable via environment variables. The current tool does not restore them.

\textbf{Credential terms.} The licensing line runs along the credential, not the
model: a Claude Pro/Max subscription is licensed for Anthropic's own clients
(the Claude apps and Claude Code), so spending it on a third-party client is the
violation --- for \emph{any} model, Claude included. In this model $subauth$ is
the subscription-OAuth mode, legitimate only on the Claude Code path. The
current tool does \emph{not} guard this: it passes $auth$ through
($vk|subscription$) without checking that a subscription credential only backs
Claude Code, so a misconfiguration could point a subscription at a non-Anthropic
harness. The target spec closes this with an enforced precondition; here it is a
known gap.

\end{document}
